\documentclass{article}
\usepackage[UTF8]{ctex}
\usepackage{amsmath}
\usepackage{booktabs}
\usepackage{graphicx}
\usepackage{listings}
\usepackage{xcolor}

\lstset{
    language=Python,
    basicstyle=\ttfamily\small,
    keywordstyle=\color{blue},
    stringstyle=\color{red},
    commentstyle=\color{green!60!black},
    numbers=left,
    numberstyle=\tiny\color{gray},
    frame=single,
    backgroundcolor=\color{gray!5},
    breaklines=true,
    captionpos=b
}

\begin{document}

\title{数学建模TOPSIS法超通俗举例+完整步骤}
\author{}
\date{}
\maketitle

TOPSIS法是\textbf{多属性决策}方法，用来给多个对象打分排名，核心：找最优解、最劣解，算各对象到两者的距离，越靠近最优、越远离最劣排名越高。

\section{举例场景}
评价3个学生：数学、英语、体育3门成绩，\textbf{越高越好}，用TOPSIS排名。

原始数据：

\begin{table}[htbp]
    \centering
    \begin{tabular}{cccc}
        \toprule
        学生 & 数学 & 英语 & 体育 \\
        \midrule
        A & 80 & 90 & 70 \\
        B & 90 & 85 & 80 \\
        C & 70 & 80 & 90 \\
        \bottomrule
    \end{tabular}
    \caption{学生成绩数据}
    \label{tab:student_scores}
\end{table}

\section{完整计算步骤}

\subsection{步骤1：数据正向化（本例全是高优，直接用）}

\subsection{步骤2：归一化（统一量纲）}

\begin{lstlisting}
import numpy as np

# 原始矩阵
x = np.array([
    [80, 90, 70],
    [90, 85, 80],
    [70, 80, 90]
])

# 每列平方和开根号
col_norm = np.sqrt(np.sum(x**2, axis=0))
# 归一化矩阵
r = x / col_norm
print("归一化矩阵：")
print(np.round(r, 4))
\end{lstlisting}

\begin{verbatim}
归一化矩阵：
[[0.5744 0.6106 0.5026]
 [0.6462 0.5767 0.5744]
 [0.5026 0.5428 0.6462]]
\end{verbatim}

\begin{lstlisting}
# 最优解（每列最大）
best = np.max(r, axis=0)
# 最劣解（每列最小）
worst = np.min(r, axis=0)

print("最优解：", np.round(best, 4))
print("最劣解：", np.round(worst, 4))
\end{lstlisting}

\begin{verbatim}
最优解： [0.6462 0.6106 0.6462]
最劣解： [0.5026 0.5428 0.5026]
\end{verbatim}

\begin{lstlisting}
# 到最优解距离
d_best = np.sqrt(np.sum((r - best)**2, axis=1))
# 到最劣解距离
d_worst = np.sqrt(np.sum((r - worst)**2, axis=1))

print("A/B/C到最优距离：", np.round(d_best, 4))
print("A/B/C到最劣距离：", np.round(d_worst, 4))
\end{lstlisting}

\begin{verbatim}
A/B/C到最优距离： [0.1605 0.0794 0.1588]
A/B/C到最劣距离： [0.0988 0.1641 0.1436]
\end{verbatim}

\subsection{步骤3：确定最优解、最劣解}
最优解：每列最大值  
最劣解：每列最小值  

\subsection{步骤4：计算各对象到最优、最劣的欧式距离}

\subsection{步骤5：计算综合得分}
公式：  
得分 $S = \frac{劣距}{优距 + 劣距}$

结果：
- A：0.381
- B：0.674
- C：0.475

\subsection{步骤6：排名}
B > C > A

\section{TOPSIS核心总结}
1. 用途：多指标综合评价、排名选优  
2. 关键：归一化+找两极+算距离+算得分  
3. 优点：简单通用、结果稳定  

\end{document}